\documentclass[11pt,a4paper]{article}
\usepackage[margin=2.2cm,top=2.0cm,bottom=2.2cm]{geometry}
\usepackage{amsmath,amssymb}
\usepackage[T1]{fontenc}
\usepackage{mathptmx}
\usepackage{xcolor}
\usepackage{enumitem}
\usepackage{fancyhdr}
\usepackage{tabularx}
\usepackage{mdframed}
\usepackage{multicol}

\definecolor{navy}{HTML}{1F4E79}
\definecolor{rule}{HTML}{333333}
\definecolor{boxbg}{HTML}{F4F6F9}
\definecolor{keyred}{HTML}{B02020}

% \ANSWERStrue for the tutor key
\newif\ifANSWERS
\input{mode}

\pagestyle{fancy}\fancyhf{}
\renewcommand{\headrulewidth}{0pt}
\renewcommand{\footrulewidth}{0.4pt}
\fancyfoot[L]{\footnotesize\color{navy}SAT Math: Foundations to Mastery \ $\cdot$\ Quadratics\ifANSWERS\ $\cdot$\ \textbf{\color{keyred}ANSWER KEY}\fi}
\fancyfoot[R]{\footnotesize\color{navy}Page \thepage\ of \pageref{LastPage}}
\usepackage{lastpage}

\setlength{\parindent}{0pt}
\setlength{\parskip}{4pt}

\newcommand{\PartHead}[2]{%
  \vspace{6pt}
  {\color{navy}\rule{\textwidth}{0.8pt}}\\[2pt]
  {\large\bfseries\color{navy}#1}\hfill{\small\itshape #2}\\[-4pt]
  {\color{navy}\rule{\textwidth}{0.4pt}}
  \vspace{4pt}
}
\newcommand{\sitting}[1]{{\footnotesize\textsc{\color{navy}#1}}\par\vspace{2pt}}
\newcommand{\spr}{\\[3pt]\hspace*{6pt}Answer: \underline{\hspace{3.2cm}}\hfill{\footnotesize\itshape(write your answer)}}
\ifANSWERS
  \newcommand{\ans}[1]{\ {\color{keyred}\bfseries #1}}
  \newcommand{\note}[1]{\par{\footnotesize\color{keyred}#1}}
\else
  \newcommand{\ans}[1]{}
  \newcommand{\note}[1]{}
\fi

% four lettered options, one per line, under the stem
\newenvironment{choices}
  {\par\nopagebreak\begin{itemize}[label={},leftmargin=16pt,itemsep=0pt,topsep=2pt,parsep=0pt,%
     beginpenalty=10000,midpenalty=10000]}
  {\end{itemize}}

\begin{document}
\label{firstpage}

\begin{center}
{\LARGE\bfseries\color{navy}Quadratics}\\[4pt]
{\color{navy}SAT Math $\cdot$ Foundations to Mastery $\cdot$ Week of 25 October}
\ifANSWERS\\[4pt]{\color{keyred}\bfseries ANSWER KEY --- tutor copy}\fi
\end{center}
\vspace{2pt}
{\color{rule}\rule{\textwidth}{1pt}}
\vspace{6pt}

\ifANSWERS\else
Name: \underline{\hspace{6cm}} \hfill Due: \underline{\hspace{3.2cm}}
\vspace{6pt}
\fi

\begin{mdframed}[backgroundcolor=boxbg,linecolor=navy,linewidth=0.6pt,innertopmargin=6pt,innerbottommargin=8pt]
{\bfseries\color{navy}How to use this sheet}\par\vspace{2pt}
This is four short sittings, not one long one. Four half-hours across the week will
do more for you than one long Saturday, and you are far more likely to finish it.

\begin{itemize}[leftmargin=14pt,itemsep=1pt,topsep=3pt]
  \item \textbf{This is new ground.} Parts A and B are built so you can finish them. Expect to go slowly in Part C, and expect that to be normal.
  \item \textbf{A squared bracket has two answers, not one.} Whenever you take a square root, write $\pm$ before you write anything else. That single habit is the point of the week.
  \item When you get one wrong, write it in your error log: what you did, and what you should have done.
\end{itemize}
\end{mdframed}

\vspace{4pt}
\begin{center}
\begin{tabularx}{\textwidth}{@{}l l X r@{}}
\textbf{Sitting} & \textbf{Part} & \textbf{What it drills} & \textbf{Minutes}\\[2pt]
\hline\rule{0pt}{11pt}
Sunday   & A    & Retention: translation setups                 & 15\\
Tuesday  & B    & Expanding and factoring                       & 25\\
Thursday & C    & Solving quadratics                            & 25\\
Saturday & D, E & In context; six mixed, timed                  & 27\\
\hline
\end{tabularx}
\end{center}

% ------------------------------------------------ PART A
\PartHead{Part A --- Retention}{No calculator. 15 minutes.}
\sitting{Sunday}
The first six are last week's work: \textbf{write the equation only, do not solve it.}
The last four keep weeks two and three alive. Three weeks running at once, in fifteen minutes.

\begin{enumerate}[leftmargin=20pt,itemsep=7pt]
\item A gym charges a \$40 joining fee plus \$29 for each month. Write an equation for the total cost $C$ after $m$ months.\ifANSWERS\else\\
  \underline{\hspace{\linewidth}}\fi
  \ans{\\ $C=40+29m$. The fee is paid once, so it is not attached to $m$.}

\item Five more than twice a number $n$ is 31. Write the equation.\ifANSWERS\else\\
  \underline{\hspace{\linewidth}}\fi
  \ans{\\ $2n+5=31$. ``Five more than twice $n$'' is $2n+5$, not $2(n+5)$.}

\item A rectangle's length $L$ is 4 cm less than three times its width $W$, and its perimeter is 72 cm. Write two equations.\ifANSWERS\else\\
  \underline{\hspace{\linewidth}}\\[4pt]
  \underline{\hspace{\linewidth}}\fi
  \ans{\\ $L=3W-4$ \ and \ $2L+2W=72$. ``4 less than three times $W$'' is $3W-4$. Writing $L+W=72$ for the perimeter is the usual slip.}

\item A shop sells 150 items in all, hats at \$6 and scarves at \$11, taking \$1{,}245. Let $h$ be the number of hats and $s$ the number of scarves. Write two equations.\ifANSWERS\else\\
  \underline{\hspace{\linewidth}}\\[4pt]
  \underline{\hspace{\linewidth}}\fi
  \ans{\\ $h+s=150$ \ and \ $6h+11s=1245$. One count equation, one money equation.}

\item A tank holds 400 litres and drains at 12 litres each minute. Write an equation for the volume $V$ left after $t$ minutes.\ifANSWERS\else\\
  \underline{\hspace{\linewidth}}\fi
  \ans{\\ $V=400-12t$. Draining makes the rate negative.}

\item The sum of three consecutive integers is 84. Let $n$ be the smallest. Write the equation.\ifANSWERS\else\\
  \underline{\hspace{\linewidth}}\fi
  \ans{\\ $n+(n+1)+(n+2)=84$, which tidies to $3n+3=84$.}

\item A line passes through $(-3,8)$ and $(5,-4)$. Its slope is \ifANSWERS\else\underline{\hspace{2cm}}\fi \ans{$\frac{-4-8}{5-(-3)}=\frac{-12}{8}=-1.5$}

\item For $4x-5y=20$, the $y$-coordinate of the $y$-intercept is \ifANSWERS\else\underline{\hspace{2cm}}\fi \ans{$-4$ \ (put $x=0$: $-5y=20$)}

\item $y=2x+1$ \ and \ $3x+y=16$. \ What is the value of $x$? \spr
  \ans{\\ $x=3$ (and $y=7$). Substitution: $3x+(2x+1)=16$, so $5x=15$. Check: $9+7=16$.}

\item $5x+2y=4$ \ and \ $3x-2y=12$. \ What is the value of $x$? \spr
  \ans{\\ $x=2$ (and $y=-3$). Elimination: $+2y$ and $-2y$ are opposites, so adding gives $8x=16$.}
\end{enumerate}

\note{\textbf{Marking rule: below 8 of 10 means last week's translation work has not stuck.} Items 3 and 4 are the diagnostic ones --- they need two equations, and a student who writes one has not seen the pattern yet.}

% ------------------------------------------------ PART B
\PartHead{Part B --- Expanding and factoring}{Calculator allowed. 25 minutes.}
\sitting{Tuesday}
This is the whole of Advanced Math in one skill: moving between $(x+a)(x+b)$ and
$x^2+bx+c$, in both directions. Expanding is multiplication you already know.
Factoring is the same thing read backwards --- find two numbers that multiply to the
last term and add to the middle one. Nothing here is beyond you.

Expand and simplify.
\begin{enumerate}[leftmargin=20pt,itemsep=7pt]
\item $(x+5)(x+3)$ \ifANSWERS\ans{$=x^2+8x+15$}\else\ \underline{\hspace{5.4cm}}\fi
\item $(x-7)(x+2)$ \ifANSWERS\ans{$=x^2-5x-14$ \ (the middle term is $2x-7x$)}\else\ \underline{\hspace{5.4cm}}\fi
\item $(2x+3)(x-5)$ \ifANSWERS\ans{$=2x^2-7x-15$ \ ($-10x+3x$)}\else\ \underline{\hspace{5.4cm}}\fi
\item $(x+6)(x-6)$ \ifANSWERS\ans{$=x^2-36$ \ (the $x$ terms cancel: difference of squares)}\else\ \underline{\hspace{5.4cm}}\fi
\item $(x-5)^2$ \ifANSWERS\ans{$=x^2-10x+25$, \emph{not} $x^2+25$}\else\ \underline{\hspace{5.4cm}}\fi
\end{enumerate}

\vspace{2pt}
Factor.
\begin{enumerate}[leftmargin=20pt,itemsep=7pt,start=6]
\item $x^2-49$ \ifANSWERS\ans{$=(x+7)(x-7)$}\else\ \underline{\hspace{5.4cm}}\fi
\item $x^2+12x+36$ \ifANSWERS\ans{$=(x+6)^2$ \ (a perfect square: $6+6=12$ and $6\times6=36$)}\else\ \underline{\hspace{5.4cm}}\fi
\item $x^2+9x+20$ \ifANSWERS\ans{$=(x+4)(x+5)$}\else\ \underline{\hspace{5.4cm}}\fi
\item $x^2-3x-28$ \ifANSWERS\ans{$=(x-7)(x+4)$ \ (multiply to $-28$, add to $-3$)}\else\ \underline{\hspace{5.4cm}}\fi
\item $2x^2-8x-24$ \ifANSWERS\ans{$=2(x-6)(x+2)$ \ (pull the 2 out first, then factor $x^2-4x-12$)}\else\ \underline{\hspace{5.4cm}}\fi
\end{enumerate}

\vspace{2pt}
Choose one answer.
\begin{enumerate}[leftmargin=20pt,itemsep=9pt,start=11]
\item Which expression is equivalent to $3x^2+15x-42$?
  \begin{choices}
  \item A) $3(x-7)(x+2)$
  \item B) $3(x+7)(x-2)$
  \item C) $3(x+14)(x-1)$
  \item D) $(x+7)(x-2)$
  \end{choices}
  \ans{\\ B. Pull out the 3 to get $3(x^2+5x-14)$, then two numbers multiplying to $-14$ and adding to $+5$ are $+7$ and $-2$. \\ A swaps the signs in the factor pair and gives $3x^2-15x-42$. C uses the wrong factor pair of $-14$ and gives $3x^2+39x-42$. D drops the common factor on the way back out.}

\item Which expression is equivalent to $(x-9)^2$?
  \begin{choices}
  \item A) $x^2-81$
  \item B) $x^2-18x-81$
  \item C) $x^2+18x+81$
  \item D) $x^2-18x+81$
  \end{choices}
  \ans{\\ D. $(x-9)(x-9)=x^2-9x-9x+81$. \\ A squares each term separately and loses the middle term entirely --- the most common error on this page. B forgets that $(-9)(-9)$ is positive. C loses the minus in the middle term.}
\end{enumerate}

% ------------------------------------------------ PART C
\PartHead{Part C --- Solving quadratics}{Calculator allowed. 25 minutes.}
\sitting{Thursday}
A quadratic almost always has \emph{two} answers. Before you start each item, write
the number 2 in the margin, and do not move on until you have found both.

Two routes. If one side is already zero, factor and set each bracket to zero. If the
equation looks like $(x-h)^2=k$, take the square root of both sides --- and write
$\pm$ on the right-hand side \emph{as you write it}, not afterwards.

\begin{enumerate}[leftmargin=20pt,itemsep=9pt]
\item $x^2-7x+12=0$. \ What is the greater of the two solutions? \spr
  \ans{\\ $4$. $(x-3)(x-4)=0$, so $x=3$ or $x=4$. Check: $16-28+12=0$.}

\item $x^2+2x-15=0$. \ What is the lesser of the two solutions? \spr
  \ans{\\ $-5$. $(x+5)(x-3)=0$, so $x=-5$ or $x=3$. Check: $25-10-15=0$.}

\item $x^2-5x=24$. \ What is the greater of the two solutions? \spr
  \ans{\\ $8$. Set it to zero first: $x^2-5x-24=0$, so $(x-8)(x+3)=0$ and $x=8$ or $x=-3$. Check: $64-40=24$. Factoring before moving the 24 across gets you nowhere.}

\item $x^2=6x+16$. \ What is the lesser of the two solutions? \spr
  \ans{\\ $-2$. Set it to zero: $x^2-6x-16=0$, so $(x-8)(x+2)=0$ and $x=8$ or $x=-2$. Check: $(-2)^2=4$ and $6(-2)+16=4$.}

\item $(x-3)^2=49$. \ What is the lesser of the two solutions?
  \begin{choices}
  \item A) $-10$
  \item B) $-4$
  \item C) $4$
  \item D) $10$
  \end{choices}
  \ans{\\ B. $x-3=\pm7$, so $x=10$ or $x=-4$. Check $x=-4$: $(-7)^2=49$. \\ D is the positive root, taken by a student who wrote $x-3=7$ and stopped. C is $-4$ with the minus dropped on the way to the page. A is $-10$, the positive root negated instead of solved for.}

\item $(x+2)^2=36$. \ What is the lesser of the two solutions?
  \begin{choices}
  \item A) $-8$
  \item B) $-4$
  \item C) $4$
  \item D) $8$
  \end{choices}
  \ans{\\ A. $x+2=\pm6$, so $x=4$ or $x=-8$. Check $x=-8$: $(-6)^2=36$. \\ C is the positive root, kept alone. D is $-8$ with the minus lost. B subtracts 2 from $-6$ with a sign slip.}

\item $(x+6)^2=25$. \ What is the lesser of the two solutions?
  \begin{choices}
  \item A) $-11$
  \item B) $-1$
  \item C) $1$
  \item D) $11$
  \end{choices}
  \ans{\\ A. $x+6=\pm5$, so $x=-1$ or $x=-11$. Check $x=-11$: $(-5)^2=25$. \\ B is the root you reach from $+5$ alone. D is $-11$ with the minus dropped. C is $-1$ with the minus dropped. Here \emph{both} solutions are negative, which catches a student who assumes one of each.}

\item $(x-1)^2=64$. \ What is the lesser of the two solutions?
  \begin{choices}
  \item A) $-9$
  \item B) $-7$
  \item C) $7$
  \item D) $9$
  \end{choices}
  \ans{\\ B. $x-1=\pm8$, so $x=9$ or $x=-7$. Check $x=-7$: $(-8)^2=64$. \\ D is the positive root alone. C is $-7$ with the minus dropped. A negates the positive root rather than solving for the second one.}

\item $2x^2-7x+3=0$. \ Use the quadratic formula. What is the greater of the two solutions? \spr
  \ans{\\ $3$. $x=\frac{7\pm\sqrt{49-24}}{4}=\frac{7\pm5}{4}$, giving $3$ and $0.5$. Check $x=3$: $18-21+3=0$. The sign trap is $-b$: with $b=-7$ the top of the fraction starts $+7$.}

\item $3x^2+5x-2=0$. \ Use the quadratic formula. What is the lesser of the two solutions? \spr
  \ans{\\ $-2$. $x=\frac{-5\pm\sqrt{25+24}}{6}=\frac{-5\pm7}{6}$, giving $\tfrac13$ and $-2$. Check $x=-2$: $12-10-2=0$. Note $-4ac=-4(3)(-2)=+24$, so the discriminant grows.}

\item How many real solutions does $x^2+4x+9=0$ have?
  \begin{choices}
  \item A) Zero
  \item B) Exactly one
  \item C) Exactly two
  \item D) Infinitely many
  \end{choices}
  \ans{\\ A. The discriminant is $b^2-4ac=16-36=-20$. A negative discriminant means no real solution, because no real number squares to a negative. \\ C is the default assumption that every quadratic has two roots. B is the answer for a discriminant of exactly zero. D is a habit carried over from systems of equations.}

\item What is the sum of the two solutions of $x^2-11x+24=0$? \spr
  \ans{\\ $11$. The solutions are 3 and 8, and $3+8=11$. Faster: for $x^2+bx+c=0$ the two roots always add to $-b$, so you can read 11 straight off without solving.}
\end{enumerate}

% ------------------------------------------------ PART D
\PartHead{Part D --- In context}{Calculator allowed. 15 minutes.}
\sitting{Saturday}
Here the algebra gives you two answers and the situation throws one away. A time
before the ball was thrown, a width of $-10$ metres: the equation does not know
these are impossible, so you have to. Solve, then look at both answers and ask which
one could happen.

\begin{enumerate}[leftmargin=20pt,itemsep=10pt]
\item A ball is thrown from a roof. Its height in feet after $t$ seconds is
  $h(t)=-16t^2+48t+64$. After how many seconds does it hit the ground? \spr
  \ans{\\ $4$. Set $h=0$: $-16(t^2-3t-4)=0$, so $(t-4)(t+1)=0$ and $t=4$ or $t=-1$. Check $t=4$: $-256+192+64=0$.}\\[2pt]
  The other solution is rejected. Why?\ifANSWERS\else\\
  \underline{\hspace{\linewidth}}\fi
  \ans{\\ $t=-1$ is one second \emph{before} the ball was thrown, and the model only describes what happens from the throw onwards. Negative time is not a moment in this story.}

\item A rectangular garden is 3 metres longer than it is wide, and its area is
  70 square metres. What is its width, in metres? \spr
  \ans{\\ $7$. $w(w+3)=70$, so $w^2+3w-70=0$ and $(w-7)(w+10)=0$, giving $w=7$ or $w=-10$. Check: $7\times10=70$.}\\[2pt]
  The other solution is rejected. Why?\ifANSWERS\else\\
  \underline{\hspace{\linewidth}}\fi
  \ans{\\ A width of $-10$ metres is not a length. A side of a rectangle has to be positive, so only $w=7$ describes a real garden.}

\item A stall's daily revenue, in dollars, is modelled by $R(p)=-5p^2+90p$, where
  $p$ is the price in dollars. For which lower of the two prices is the revenue
  exactly \$400? \spr
  \ans{\\ $8$. $-5p^2+90p=400$ gives $p^2-18p+80=0$, so $(p-8)(p-10)=0$ and $p=8$ or $p=10$. Check $p=8$: $-320+720=400$. Here both answers are possible prices, so nothing is rejected --- that is the point of putting it next to the others.}

\ifANSWERS\else\newpage\fi
\item A ball is thrown upward and its height in feet after $t$ seconds is
  $h(t)=-16t^2+64t+5$. What is the greatest height, in feet, that the ball reaches?
  \begin{choices}
  \item A) $2$
  \item B) $5$
  \item C) $64$
  \item D) $69$
  \end{choices}
  \ans{\\ D. The vertex is at $t=\frac{-64}{2(-16)}=2$, and $h(2)=-64+128+5=69$. \\ A is $t=2$, the time rather than the height --- the question asked how high, not when. B is the starting height at $t=0$. C reads the 64 out of the equation without doing anything with it.}

\item A stall's daily profit, in dollars, is $P(n)=n^2+10n-1200$, where $n$ is the
  number of items sold. How many items must it sell to break even, that is, for the
  profit to be \$0? \spr
  \ans{\\ $30$. $n^2+10n-1200=0$ has discriminant $100+4800=4900=70^2$, so $n=\frac{-10\pm70}{2}$, giving $30$ and $-40$. Check: $900+300-1200=0$.}\\[2pt]
  The other solution is rejected. Why?\ifANSWERS\else\\
  \underline{\hspace{\linewidth}}\fi
  \ans{\\ $n=-40$ would mean selling a negative number of items, which cannot happen. A count of things is never negative.}
\end{enumerate}

% ------------------------------------------------ PART E
\PartHead{Part E --- Mixed, timed}{Calculator allowed. 12 minutes, timed.}
\sitting{Saturday}
Set a timer for 12 minutes. These span Parts B to D. Answer every one --- there is no
penalty for a wrong answer, so a blank is a mark thrown away. Where you find two
solutions, read the question again before you write one of them down.

\begin{enumerate}[leftmargin=20pt,itemsep=9pt]
\item Which expression is equivalent to $(x-6)(x+2)$?
  \begin{choices}
  \item A) $x^2-12$
  \item B) $x^2-4x-12$
  \item C) $x^2+4x-12$
  \item D) $x^2-8x-12$
  \end{choices}
  \ans{\\ B. The middle term is $2x-6x=-4x$. \\ A multiplies the first and last terms only. C has the sign of the middle term backwards. D adds $-6$ and $-2$ instead of $-6$ and $+2$.}

\item $x^2-9x+18=0$. \ What is the greater of the two solutions? \spr
  \ans{\\ $6$. $(x-3)(x-6)=0$. Check: $36-54+18=0$.}

\item $(x+5)^2=81$. \ What is the lesser of the two solutions?
  \begin{choices}
  \item A) $-14$
  \item B) $-4$
  \item C) $4$
  \item D) $14$
  \end{choices}
  \ans{\\ A. $x+5=\pm9$, so $x=4$ or $x=-14$. Check $x=-14$: $(-9)^2=81$. \\ C is the positive root kept alone. D is $-14$ with the minus dropped. B comes from $-9+5$ with a sign slip.}

\item Which expression is equivalent to $4x^2-25$?
  \begin{choices}
  \item A) $(2x-5)(2x+5)$
  \item B) $(2x-5)^2$
  \item C) $(4x-5)(x+5)$
  \item D) $(4x-25)(x+1)$
  \end{choices}
  \ans{\\ A. Both terms are squares, $4x^2=(2x)^2$ and $25=5^2$, so this is a difference of squares. \\ B is the perfect square confused with the difference of squares, and expands to $4x^2-20x+25$. C and D split the 4 and the 25 into a pair that does not cancel the middle term: they give $4x^2+15x-25$ and $4x^2-21x-25$.}

\item A ball is thrown from a ledge and its height in feet after $t$ seconds is
  $h(t)=-16t^2+32t+48$. After how many seconds does it hit the ground? \spr
  \ans{\\ $3$. $-16(t^2-2t-3)=0$, so $(t-3)(t+1)=0$ and $t=3$ or $t=-1$. Reject $t=-1$: the clock starts at the throw. Check $t=3$: $-144+96+48=0$.}

\item What is the sum of the two solutions of $2x^2-14x+5=0$? \spr
  \ans{\\ $7$. The roots are not whole numbers, but their sum is $-\frac{b}{a}=\frac{14}{2}=7$. The discriminant is $196-40=156$, so there are two real roots and you never need to write either one down.}
\end{enumerate}

\ifANSWERS\else
\newpage
\PartHead{Error log}{Fill this in as you go.}
\begin{center}
\begin{tabularx}{\textwidth}{|c|c|X|X|}
\hline
\textbf{Date} & \textbf{Q} & \textbf{What I did} & \textbf{What I should have done}\\
\hline
\rule{0pt}{22pt} & & &\\ \hline
\rule{0pt}{22pt} & & &\\ \hline
\rule{0pt}{22pt} & & &\\ \hline
\rule{0pt}{22pt} & & &\\ \hline
\rule{0pt}{22pt} & & &\\ \hline
\rule{0pt}{22pt} & & &\\ \hline
\rule{0pt}{22pt} & & &\\ \hline
\end{tabularx}
\end{center}
\vspace{4pt}
\begin{center}\small\itshape
Bring this sheet and your error log to the next session.
\end{center}
\fi

\ifANSWERS
\vspace{8pt}
\PartHead{Five-minute marking routine}{Tutor}
\begin{enumerate}[leftmargin=20pt,itemsep=4pt]
\item \textbf{Score Part A only.} Below 8/10 means last week's translation work has not stuck, and the next session opens by repairing that rather than pushing further into quadratics.
\item \textbf{In Part C, count how many times he gave only the positive root.} Items 5 to 8 are built for it. That single error is the one to name out loud --- give it a number, say it once, and let him hear that every other mistake on the page is smaller than this one.
\item \textbf{Pick two wrong answers from Part E} and have him narrate his thinking out loud. Leave the rest unmarked.
\end{enumerate}
This is his first week of Advanced Math and he answered nothing in it on the diagnostic. Any part finished is progress worth saying out loud.
\fi

\end{document}
